\section{Physical Real2Sim: el umbral real de una buena simulación}

Tras explicar por qué la inteligencia corporizada debe apoyarse en la simulación, esta sección entra en la cuestión técnica más difícil: qué cuenta como una buena simulación. Xie Chen (谢晨) distingue entre el Real2Sim visual y el Real2Sim físico. El Real2Sim visual reconstruye la apariencia, la geometría y la escena tal como se ven; el Real2Sim físico, en cambio, debe trasladar a la simulación las propiedades mecánicas, las relaciones de interacción, el cuerpo del robot y los parámetros físicos del mundo real. Esta segunda versión es la que la inteligencia corporizada realmente necesita.

Una buena simulación requiere al menos cuatro componentes: activos y escenas sim-ready de alta calidad, un solver físico preciso, renderizado y simulación de sensores, y API, metadata y una cadena de herramientas que los algoritmos de entrenamiento puedan invocar. Si falta cualquiera de ellos, surgen problemas. Por ejemplo, si la puerta de un refrigerador se ve muy realista pero no se puede abrir, o si se abre pero con una fuerza incorrecta, la política que el robot aprenda en la simulación podría no transferirse.

\begin{figure}[H]
\centering
\includegraphics[width=0.96\textwidth]{real2sim-workflow.png}
\caption{Flujo de trabajo de Real2Sim: de la captura del mundo real a la reconstrucción, la calibración y una simulación reproducible. Diagrama conceptual propio, elaborado a partir de la conversación entre 00:32:41 y 00:46:18.}
\end{figure}

\begin{knowledgebox}{Lectura de la figura: Real2Sim convierte el mundo físico en un mundo entrenable}
Lo más importante de la figura es la «calibración», no la «reconstrucción». La captura real recupera la apariencia, la geometría, los parámetros mecánicos y las trayectorias de interacción; los activos de simulación deben poder ser tocados, jalados, empujados y sujetados por el robot; por último, hay que verificar mediante la evaluación en tareas reales si la política se transfiere. Solo al recorrer la cadena completa, Real2Sim deja de ser modelado 3D y se convierte en infraestructura de entrenamiento.
\end{knowledgebox}

\subsection{Cuatro criterios de una buena simulación}

Esta subsección concreta qué es una «buena simulación». Primero, los activos y las escenas deben contener suficiente información física, incluidos cuerpos de colisión, articulaciones, fricción, peso, par y materiales. Segundo, el solver subyacente debe ser lo bastante preciso y capaz de resolver problemas distintos, como cuerpos rígidos, cuerpos no rígidos, telas y líquidos. Tercero, la salida de la simulación debe servir al entrenamiento de algoritmos mediante API y metadata, no solo mostrarse a personas. Cuarto, la simulación debe ser lo bastante eficiente para dar servicio al aprendizaje por refuerzo en paralelo, sobre todo al RL con visión en el lazo.

Hay aquí un punto que se pasa por alto con facilidad: el cuerpo del robot en la simulación también tiene que ser preciso. En la entrevista se menciona que la mano del robot real de un cliente podía cargar uno o dos kilogramos, pero en la simulación solo podía cargar 0.1 kilogramos; el equipo pasó seis meses sin lograr ajustarlo. Esto muestra que la Sim2Real gap no proviene solo de la escena, sino también del propio modelo del robot. Una mano errónea, un límite articular erróneo o una curva de par errónea bastan para distorsionar los resultados del entrenamiento.

\begin{importantbox}{Criterios de una buena simulación}
Una buena simulación no es la que tiene buena imagen, sino la que permite llevar a robots reales los algoritmos entrenados en ella. Debe contar con parámetros físicos reales, un solver preciso, una cadena de herramientas entrenable, alta capacidad de paralelización y un ciclo de evaluación que verifique continuamente el éxito o el fracaso de Sim2Real.
\end{importantbox}

\begin{knowledgebox}{Términos clave: Real2Sim físico}
\begin{tabularx}{\textwidth}{p{0.20\textwidth}p{0.35\textwidth}X}
\toprule
Término & Significado & Por qué afecta el entrenamiento \\
\midrule
Solver & Resolvedor del motor físico que calcula el movimiento, las colisiones y las restricciones & Determina si los resultados mecánicos son confiables. \\
Collision Body & Geometría que participa en el cálculo de colisiones & Sin cuerpos de colisión correctos, sujetar objetos y abrir puertas deja de ser realista. \\
Sim-ready Asset & Activo que puede usarse directamente en interacciones de simulación & No basta con la apariencia del modelo; también debe incluir parámetros físicos e interfaces. \\
Metadata & Etiquetas, estados, parámetros e información semántica adicionales & El entrenamiento y la evaluación necesitan leer esta información estructurada. \\
Parallel Environment & Múltiples instancias de simulación que se ejecutan en paralelo en RL & Determina el rendimiento del entrenamiento y la eficiencia de exploración. \\
\bottomrule
\end{tabularx}
\end{knowledgebox}

\subsection{Diferencias entre el aprendizaje por imitación y el aprendizaje por refuerzo}

Esta subsección examina los requisitos de la simulación desde el objetivo de entrenamiento. El aprendizaje por imitación puede usar datos de lazo abierto: una persona o un sistema de teleoperación demuestra una trayectoria y el modelo aprende a replicarla. Es adecuado para demos tempranas, pero su generalización suele ser limitada. Por ejemplo, el modelo puede aprender a abrir la tapa de cierta botella alta, pero no necesariamente generaliza a botellas bajas, a otras formas de girar o a otros materiales. El aprendizaje por refuerzo, en cambio, necesita un entorno de lazo cerrado en el que el robot pruebe y se equivoque una y otra vez y aprenda de la reward, lo que exige que la simulación sea a la vez precisa y eficiente.

El ponente compara el RL de robots con el posentrenamiento de los modelos grandes: primero se tiene un modelo base y luego se hace fine-tune mediante aprendizaje por refuerzo. La inteligencia corporizada también podría necesitar un VLA o un modelo base robótico, cuya política se ajuste después mediante RL en una simulación de alta calidad. La dificultad es que los entornos de RL con visión en el lazo son muy pesados: cuántos environment pueden ejecutarse en paralelo en una GPU y qué tan alto es el total FPS determinan directamente si el entrenamiento puede escalar.

\begin{warningbox}{Si la simulación no admite RL, difícilmente sostendrá la generalización de la siguiente etapa}
Una simulación que solo sirve al aprendizaje por imitación puede generar trayectorias y demostraciones; una simulación que sirve al RL debe permitir que la política actúe, observe, falle, reintente y obtenga reward. Esta última exige mucho más en precisión física, eficiencia en paralelo y cadena de herramientas.
\end{warningbox}

\subsection{Simulación con calidad de agua potable}

Esta subsección resume la visión de Guanglun Zhineng (光轮智能): una buena simulación debería ser tan accesible como el agua potable. La metáfora es contundente, porque convierte la simulación de «herramienta misteriosa de algún laboratorio» en infraestructura pública de la industria. Hoy muchas simulaciones son como agua sin filtrar: parece haber mucha, pero beberla causa problemas; lo que de verdad impulsará a la industria en el futuro es que todo investigador en robótica pueda acceder a entornos de simulación de calidad suficiente, suficientemente utilizables y suficientemente verificables.

Para lograrlo, no basta con construir un simulador. En la ronda final de preguntas rápidas, Xie Chen subraya que «simulación no equivale a simulador»: el simulador es solo la capa del motor físico, mientras que la simulación también abarca activos, escenas, renderizado, API, framework, metadata, control de calidad y evaluación. Muchas empresas empiezan directamente por el simulador de bajo nivel y, en consecuencia, pueden pasar por alto los activos y las escenas, que son lo que realmente escasea. El camino de Guanglun es construir primero activos, escenas y cadenas de herramientas de alta calidad, y solo al final profundizar en el simulador de bajo nivel.

\begin{importantbox}{Simulación no equivale a simulador}
El simulador es una condición necesaria pero no suficiente. Lo que realmente le falta a la inteligencia corporizada son activos físicos interactivos, escenas escalables, interfaces entrenables, estándares de calidad verificables y un ciclo cerrado capaz de llevar de vuelta al sistema los fallos reales.
\end{importantbox}

\subsection{Resumen de la sección}

Physical Real2Sim es el núcleo técnico de EP109. Exige que el equipo lleve a la simulación la información física del mundo real con bajo costo y a escala, y que la calibre continuamente con los resultados del entrenamiento de modelos y del despliegue real. El criterio de una buena simulación no es el realismo visual, sino si la política entrenada funciona en el mundo real.
